\hypertarget{the-settled-policy}{%
\chapter{13. The settled policy}\label{the-settled-policy}}

The policy from the previous chapter is what an operator writes. It is
not what enforces. Between the two is a step skipped over until now: the
written policy (with its inheritance chain, its deltas, its includes,
and the signatures on each) has to be resolved into something a kennel
can be held to, and that resolution is a great deal of fallible,
security-critical work to be doing at the moment a workload starts. So
it is not done then. The authored artefact and the enforced artefact are
different things, and the work of turning one into the other happens
once, in advance, at a step called compilation.

This is not novel; it is how the systems Kennel is measured against
already work. AppArmor authors a text profile and compiles it into a
binary the kernel loads, and the text is never consulted again at
enforcement. SELinux compiles a monolithic binary policy from its source
modules. The authored form is for people to read and reason about; the
enforced form is for the machine to apply; and the compile step between
them is where the expensive and dangerous parts are paid for,
deliberately and once.

\hypertarget{what-compilation-produces}{%
\section{13.1 What compilation
produces}\label{what-compilation-produces}}

Compilation takes a leaf policy and produces a settled policy: a single
flat document in which the inheritance chain has been folded down, the
includes merged, the deltas applied, and every source signature and pin
verified. It carries no template reference, no include, no delta
operator: nothing left to resolve. Only the final effective rules
remain, and they are signed as a unit by whoever compiled them.

Everything hard happens here. The chain is walked from the leaf to the
root and every link's signature checked against the tier entitled to
have signed it; the includes are merged and any conflict between them is
refused rather than guessed; the deltas are applied; the framework and
template invariants are validated; the threat tags are checked against
the catalogue; and the parts of the policy fixed for the whole
installation are filled in. The result records its own provenance (every
input that went into it, named by the key that signed it), so the
settled policy can be audited later without the templates it came from
being present at all. What a kennel starts from, then, is not a policy
to resolve but a policy already resolved: read, checked once, and
frozen.

Who compiles a settled policy is not fixed to one tier. The common path
has the operator compile their own leaf, but nothing obliges the tiers
above to ship only adaptable templates. A maintainer or host can author
a complete confinement, compile it, and publish it settled, ready to run
with no leaf to write, beside the templates it offers for adaptation.
The operator then has two shapes to choose between: inherit a template
and write a leaf that adapts it, or run a finished policy a tier already
baked and signed. Both are settled policies the runtime accepts on one
signature, and either way the floor under the result is the framework's
to re-assert at spawn, so running a policy baked upstream is no more an
act of faith than running one's own.

\hypertarget{the-signature-and-the-floor}{%
\section{13.2 The signature and the
floor}\label{the-signature-and-the-floor}}

A settled policy exists to be trusted on the strength of a single
signature, and that is the larger half of what it buys. Handed one, the
runtime does not re-walk the inheritance chain, re-apply the deltas,
re-merge the includes, or re-verify the many source signatures that went
into it. It verifies one signature (the compiling authority's, over the
whole flat artefact) and on the strength of that accepts the entire body
of the policy as resolved correctly. All the fallible work of resolution
was done once, at compile; the signature is the runtime's reason to
believe it was done right. The runtime does not re-argue the policy. It
checks who vouched for it.

There is exactly one thing the signature is not trusted for, and it is
the floor. A signature says a key the system trusts vouched for this
artefact; it does not say the artefact is a valid kennel, because those
are different claims. The structural guarantees that make a kennel a
kennel (that privilege cannot escalate through it, that the categorical
denials hold, that the home directory is shadowed, that the metadata
endpoint is unreachable) are the framework's guarantees, not the
signer's. A trusted key vouching for a policy is no reason to believe
the policy upholds them, because the key's authority is over what the
operator may grant, not over what the framework must always deny.

So the framework re-asserts its own invariants directly, at the moment
of every spawn, after the one signature is verified, regardless of the
signature, and regardless of any claim the artefact carries about having
been checked already. This is the single, deliberate exception to an
otherwise complete reliance on the signature: the runtime trusts the
signer for the whole of the policy except the part the signer was never
entitled to weaken. The check is cheap, a handful of structural
assertions, and it makes one guarantee unconditional: no policy, however
it was produced and whatever key signed it, can disable the protections
that define a kennel (\texttt{reference-monitor}). A compromised
compiling authority, a mistaken template, a forged claim of validity:
none reaches the workload, because the last thing checked before the
workload starts is not the policy's provenance but the policy's floor.

\hypertarget{fail-closed}{%
\section{13.3 Fail-closed}\label{fail-closed}}

Everywhere along this path, the response to something not checking out
is the same: stop, and run nothing. A policy that does not parse is
rejected before any signature is examined: there is no verifying the
integrity of bytes that have no structure, so a malformed artefact never
reaches the part of the system that would trust a signature on it. A
signature that does not verify, an include that conflicts with another,
a reference whose bytes have changed since they were pinned, an
invariant a delta tried to weaken, a placeholder left unfilled at spawn:
each refuses the kennel outright. No workload starts on a policy that
failed a check, and no partial kennel is built from one that failed
halfway; the failure is categorical, and it names what failed.

There is no degraded mode, and that absence is deliberate. A confinement
that falls back to something weaker when it cannot apply what was asked
is a confinement an attacker works to make fail. A kennel that cannot be
resolved to a valid settled policy, or whose settled policy cannot be
verified and floored at run, does not come up at all
(\texttt{refuse-to-start}). The safe state when something is wrong is
the empty one (nothing running), and the system holds to it rather than
start a workload under a confinement it is unsure of: fail-safe defaults
{[}SaltzerSchroeder{]}, applied to starting and not only to granting.

One more property keeps the enforced artefact honest. A few grants are
derived rather than written (a writable path implies a readable one, and
so on), but the derivation is done at compile and written into the
settled policy explicitly, never left to happen invisibly at run. What a
kennel enforces is exactly what its settled policy says, down to the
grants the operator never typed, so the enforced reality can always be
reconciled with the authored intent by reading the artefact. Nothing the
operator did not intend appears, and nothing they did intend is silently
widened.

\hypertarget{authored-in-one-place-enforced-in-another}{%
\section{13.4 Authored in one place, enforced in
another}\label{authored-in-one-place-enforced-in-another}}

Putting all the resolution at compile time does more than keep the spawn
fast; it shrinks what the running system must be trusted to do. The path
that starts a kennel links none of the machinery that resolves a policy:
no template parsing, no chain-walking, no include merging, no delta
application, no verifying of source signatures. All of that crossed its
gate at compile and need never run again. What the spawn path does is
verify one signature, re-assert the floor, fill in the per-instance
values, and build the kennel: a small, fixed surface that does not grow
with the complexity of the policy it enforces.

This is what lets confinement be authored in one place and enforced in
another. An organisation can compile its policies centrally (where the
templates, the keys, and the review live) and ship only the signed
settled artefacts to the machines that run workloads. Those machines
need not hold a single template or run the resolver; their entire
policy-trust surface is one signature check and the floor re-assertion.
The operator who writes a leaf on their own machine and the fleet that
enforces a settled policy pushed from a central repository are doing the
same thing at different scales, and in both the dangerous part was paid
for before the workload ever started.

The life of a policy, then, is a sequence of refusals as much as a
sequence of steps. It is authored as intent, compiled once with every
check applied and any failure stopping the whole, settled into a flat
artefact that carries its own provenance, and run behind a single
signature and a floor the framework re-asserts for itself, trusting no
signer to have left it intact. What this leaves unanswered is the
question the signatures keep raising: whose keys does the system trust,
for what, and why is that the right division? Which tier may vouch for a
floor, which may only grant within one, and which may sign nothing that
others must rely on: that is the consent model the whole design rests
on, and it is the last thing to set down.
